% !TeX root = ../../main.tex
\subsection{Thiết kế quy trình ETL}

\subsubsection{Tổng quan quy trình ETL}
\CacBuoc
\begin{enumerate}
\item Xác định các khu vực dữ liệu: nguồn, staging và kho dữ liệu.
\item Chọn công cụ ETL và cách thực hiện (ETL hoặc ELT).
\item Xác định thứ tự chạy các bước và phần phụ thuộc giữa các bảng.
\item Vẽ sơ đồ luồng ETL.
\end{enumerate}
\TaiLieu{Ch1 tr.15--16; Ch2 tr.6, 8--9, 18.}
\KetQua{Sơ đồ và mô tả luồng ETL.}
\HinhCho{Sơ đồ luồng ETL}{fig:etl}

\subsubsection{Extract -- Trích xuất dữ liệu}
\CacBuoc
\begin{enumerate}
\item Lập danh sách nguồn dựa trên Bảng~\ref{tab:nguon}.
\item Xác định định dạng, bảng mã và kiểu dữ liệu của từng tập tin.
\item Thiết kế các bảng staging chứa dữ liệu thô.
\item Kiểm tra số dòng sau khi trích xuất so với số dòng của tập nguồn.
\end{enumerate}
\TaiLieu{Ch2 tr.8.}
\KetQua{Bảng ánh xạ từ tập nguồn đến bảng staging.}
\PhanNhom{thiết kế bước Extract}

\subsubsection{Transform -- Biến đổi dữ liệu}
\CacBuoc
\begin{enumerate}
\item Làm sạch: lập quy tắc xử lý cho từng vấn đề trong Bảng~\ref{tab:chatluong}, gồm giữ lại, sửa, gán giá trị mặc định hoặc loại bỏ.
\item Ánh xạ: chuẩn hóa kiểu dữ liệu, định dạng ngày, tên cột, tên thành phố và tên ngành hàng.
\item Tích hợp: kết nối các nguồn và xử lý mã khách hàng, dữ liệu địa lý như đã nêu ở mục~\ref{sec:quanhe}.
\item Áp dụng quy tắc nghiệp vụ: thống nhất định nghĩa doanh thu, đơn hợp lệ và giao trễ theo Bảng~\ref{tab:kpi}.
\item Tính các trường dẫn xuất và các nhóm giá trị nếu cần.
\item Xác định cách phân bổ dữ liệu giữa mức đơn hàng và mức dòng sản phẩm.
\end{enumerate}
\TaiLieu{Ch2 tr.8; Ch2-cont tr.10--13.}
\KetQua{Bảng quy tắc biến đổi và bảng ánh xạ từ cột nguồn đến cột đích.}
\PhanNhom{quy tắc biến đổi dữ liệu}

\subsubsection{Load -- Nạp dữ liệu}
\CacBuoc
\begin{enumerate}
\item Xác định thứ tự nạp giữa bảng chiều và bảng sự kiện.
\item Chọn cách nạp: nạp toàn bộ hoặc nạp tăng dần, và nêu lý do.
\item Xác định cách tra cứu khóa thay thế khi nạp bảng sự kiện.
\item Xác định cách chạy lại để không làm trùng lặp dữ liệu.
\item Xác định chỉ mục hoặc phân vùng cần tạo.
\end{enumerate}
\TaiLieu{Ch2 tr.3, 8.}
\KetQua{Thứ tự và chiến lược nạp dữ liệu.}
\PhanNhom{chiến lược nạp dữ liệu}

\subsubsection{Kiểm tra chất lượng và đối soát số liệu}
\CacBuoc
\begin{enumerate}
\item Đối soát số dòng giữa nguồn, staging và kho dữ liệu.
\item Đối soát tổng doanh thu, phí vận chuyển và số đơn hàng.
\item Kiểm tra toàn vẹn tham chiếu giữa bảng sự kiện và bảng chiều.
\item Ghi nhận số bản ghi bị loại hoặc bị sửa cùng lý do.
\end{enumerate}
\TaiLieu{Ch1 tr.9, 13, 32; Ch2-cont tr.17.}
\KetQua{Bảng kết quả đối soát sau khi chạy ETL.}
\PhanNhom{kết quả kiểm tra chất lượng}

\subsubsection{Quản lý metadata}
\CacBuoc
\begin{enumerate}
\item Lập từ điển dữ liệu cho các bảng trong kho.
\item Ghi lại nguồn gốc dữ liệu và các phép biến đổi.
\item Thiết kế bảng nhật ký cho mỗi lần chạy ETL.
\end{enumerate}
\TaiLieu{Ch2-cont tr.14--15; Ch2 tr.13.}
\KetQua{Từ điển dữ liệu và cấu trúc bảng nhật ký ETL.}
\PhanNhom{metadata của hệ thống}
